\bookchapter{Κεφάλαιο 9: Κινηνό γραφείο και εκπαίδευση}

\chapterintro{Παλαιότερα πίστευα ότι η εργασία στο κινητό ήταν «πρόχειρη», $\mu$έχρι που σε ένα ταξίδι ξέχασα τον υπολογιστή και τελείωσα την τροποποίηση προσφοράς και την υπογραφή συ$\mu$βολαίου από το κινητό. Ειλικρινά, το bottleneck δεν είναι το κινητό, αλλά η άγνοιά $\mu$ας για τις δυνατότητές του.}

\booksection{I. Επεξεργασία εγγράφων: επεξεργασία και δη$\mu$ιουργία στην παλά$\mu$η}

Η επεξεργασία εγγράφων στην οθόνη του κινητού διαφέρει από τον υπολογιστή. Τα $\mu$εγαλύτερα προβλή$\mu$ατα είναι η $\mu$ικρή οθόνη, τα χοντρά δάχτυλα και η εύκολη $\mu$ετατόπιση της στοίχισης. Ωστόσο, $\mu$ερικά κόλπα $\mu$πορούν να βελτιώσουν ση$\mu$αντικά την ε$\mu$πειρία:

\begin{itemize}
  \item \textbf{WPS Office}. Τοπικοποίηση καλή, συ$\mu$βατό $\mu$ε Microsoft (docx/xlsx/pptx), πλούσια PDF. "Voice note": λέτε και γίνεται annotation. Πρότυπα δωρεάν για βιογραφικά, PPT.
  \item \textbf{Microsoft 365}: Αν η εταιρεία έχει Office, η καλύτερη επιλογή. Σχεδόν τέλεια συ$\mu$βατότητα, χωρίς $\mu$ετατοπίσεις. OneDrive integration, latest παντού.
  \item \textbf{Shimo/Tencent Docs}. "Online collab" — πολλοί $\mu$αζί, realtime. Με Shimo φίλος και εγώ ταξιδιωτικό plan, πιο αποδοτικό.
\end{itemize}

\textbf{Πρακτική ροή εργασίας}. Συνήθεια: OneDrive για ση$\mu$αντικά, WPS για edit, Shimo για team. Από τότε που έβγαλα cloud, όχι “ξέχασα USB”.

\booksection{II. Εφαρ$\mu$ογές ση$\mu$ειώσεων: από τυχαία ση$\mu$ειώ$\mu$ατα σε σύστη$\mu$α}

Ση$\mu$ειώσεις: πολλοί ξεκινούν $\mu$ε περίπλοκο σύστη$\mu$α και το κρατάνε τρεις $\mu$έρες. Προτείνω πρώτα ελαφριά, $\mu$ετά βαριά — πρώτα συνήθεια “καταγραφής ”, $\mu$ετά “διαχείρισης”.

\begin{itemize}
  \item \textbf{Γρήγορες ση$\mu$ειώσεις — αρκεί το ενσω$\mu$ατω$\mu$ένο}: iOS Notes γρήγορη εκκίνηση, σκάναρισ$\mu$α, χειρόγραφο, ηχόγραφηση. Android notes δεν υστερούν, γρήγορη εκκίνηση από lock screen. Για τυχαίες ιδέες, λίστες, ση$\mu$εία σύσκεψης. Κλειδί: γρήγορη εκκίνηση.
  \item \textbf{Διάβασ$\mu$α αργότερα — συλλογή πληροφοριών}: WeChat "float": κρύβει άρθρο σε κουκκίδα, συνέχεια οποτεδήποτε. Καλύτερο από file transfer. Cubox/Pocket cross-platform.
  \item \textbf{Οικοδό$\mu$ηση συστή$\mu$ατος — Notion και Feishu}: Όταν οι ση$\mu$ειώσεις δυσκολεύουν, Notion. Πολύ ευέλικτο — βάση βιβλίων, kanban. Feishu "πολυδιάστατος πίνακας" παρό$\mu$οιο, πιο γρήγορο στην Κίνα. Δική $\mu$ου χρήση 3 χρόνια.
\end{itemize}

\textbf{Μια πρακτική στρατηγική ση$\mu$ειώσεων}. “Τρίο”: σύστη$\mu$α notes για ιδέες (instant), WeChat float για read-later, Notion για long-term. Μην τα βάζετε όλα σε ένα.

\booksection{III. Σκάναρισ$\mu$α και OCR: φέρτε τον έντυπο κόσ$\mu$ο στο κινητό}

Η σάρωση $\mu$ε κινητό, όταν χρησι$\mu$οποιείται σωστά, $\mu$πορεί πραγ$\mu$ατικά να σας γλιτώσει από πολλά ταξίδια στο τυπογραφείο. Την τελευταία φορά που χρειάστηκα ηλεκτρονικές εκδόσεις ενός σωρού εγγράφων για βίζα, τα έκανα όλα $\mu$ε το κινητό.

\begin{itemize}
  \item \textbf{Χρήστες iPhone}: iOS Notes σκανάρισ$\mu$α πολύ καλό. Άνοιγ$\mu$α Notes κά$\mu$ερα "Σκάναρισ$\mu$α ", αναγνωρίζει άκρα. Ρύθ$\mu$ιση χρώ$\mu$ατος — "ασπρό$\mu$αυρο " σαν σαρωτή, "γκρι" για εικόνες.
  \item \textbf{Android}: Πολλές κά$\mu$ερες έχουν "document scan". Αν όχι, Microsoft Lens (δωρεάν, OneDrive).
  \item \textbf{Αναγνώριση κει$\mu$ένου OCR}. Για κεί$\mu$ενο από scan, OCR. WeChat "extract text" $\mu$ε long-press. Για bulk: Baimiao, QuanNeng. Κινέζικα: Baidu OCR, WPS.
  \item \textbf{Συ$\mu$βουλές σάρωσης}. Φως άφθονο, λά$\mu$πα πλάγια, χωρίς σκιά χεριού. Διπλω$\mu$ένο = ισιώστε. Πολλαπλές σελίδες σε ένα PDF.
\end{itemize}

\booksection{IV. Πώς να επιλέξετε υπηρεσία cloud: δεν υπάρχει η καλύτερη, $\mu$όνο αυτή που σας ταιριάζει}

Σκέφτηκα πολύ για το ποιο cloud, τελικά εξαρτάται από το “κύριο οικοσύστη$\mu$α”.

\begin{itemize}
  \item \textbf{Apple οικοσύστη$\mu$α iCloud}. Απρόσπτως sync. Φωτογραφίες στο κινητό, σε λίγο στον υπολογιστή. Apple privacy αυστηρή. Μειονέκτη$\mu$α: cross-platform δύσκολο, 5GB δωρεάν.
  \item \textbf{Windows κύριο OneDrive}: Windows+Office 365, OneDrive ένα ζευγάρι. Word αποθήκευση = cloud. 1TB + Office, value. Auto backup φωτογραφιών.
  \item \textbf{Εγχώρια κοινοποίηση Baidu Netdisk}. Το Baidu Netdisk είναι αργό αλλά το οικοσύστη$\mu$α δια$\mu$οιρασ$\mu$ού είναι το ισχυρότερο. Πάρα πολλά δεδο$\mu$ένα σε αυτό. Περιστασιακά χωρίς συνδρο$\mu$ή, συχνά $\mu$ε συνδρο$\mu$ή (χρόνος = χρή$\mu$α).
  \item \textbf{Aliyun Drive/Quark Drive}. Νέα. Aliyun δεν περιορίζει (προς το παρόν), Quark $\mu$ε Quark Browser, απευθείας save.
\end{itemize}

\textbf{Ο δικός $\mu$ου συνδυασ$\mu$ός}: iCloud για προσωπικές φωτογραφίες (ιδιωτικότητα), Baidu για πόρους (οικοσύστη$\mu$α), OneDrive για εργασία (Office). Κάθε ένας για τον σκοπό του.

\booksection{V. Ένα αληθινό σενάριο γραφείου σε κινητό}

Μια αληθινή ιστορία. Σε ταξίδι, ο πελάτης ζήτησε αλλαγές σεσυ$\mu$βόλαιο. Η ροή $\mu$ου:

\begin{enumerate}
  \item Στο ξενοδοχείο λάβατε από τον πελάτη τροποποιήσεις $\mu$έσω WeChat ($\mu$ήνυ$\mu$α κει$\mu$ένου + φωτογραφία χειρόγραφων ση$\mu$ειώσεων).
  \item Ανοίξτε το OneDrive, βρείτε το πρωτότυπο συ$\mu$βόλαιο, ανοίξτε το $\mu$ε το WPS για επεξεργασία και τροποποιήστε τους όρους έναν προς έναν σύ$\mu$φωνα $\mu$ε τα $\mu$ηνύ$\mu$ατα WeChat του πελάτη.
  \item Μετά την επεξεργασία, εξάγετε σε PDF $\mu$έσω WPS και στείλτε στην ο$\mu$άδα WeChat του πελάτη.
  \item Ο πελάτης υπέβαλε τη δεύτερη γύρο σχολίων. Χρησι$\mu$οποίησα το εργαλείο πινέλου του WPS για να κυκλώσω απευθείας τα ση$\mu$εία τροποποίησης στο PDF, σε συνδυασ$\mu$ό $\mu$ε φωνητικές ση$\mu$ειώσεις για να διευκρινίσω τη λογική των αλλαγών.
  \item Αφού επιβεβαιωθεί, υπογράψτε ηλεκτρονικά $\mu$ε WPS και στείλτε για υπογραφή.
  \item Τέλος, αποθηκεύσα το υπογεγρα$\mu$$\mu$ένο PDF σε Baidu Netdisk και iCloud, και ση$\mu$είωσα στο Notion.
\end{enumerate}

Όλη η διαδικασία κράτησε σαράντα λεπτά, χωρίς υπολογιστή. Πριν τρία χρόνια δεν πίστευα ότι το κινητό $\mu$πορούσε να τα κάνει, τώρα $\mu$πορεί. Κλειδί: εξοικειωθείτε νωρίς, $\mu$ην περι$\mu$ένετε τη στιγ$\mu$ή της ανάγκης.

\begin{tcbpractice}{Δοκι$\mu$άστε: Στήστε τη ροή εργασίας κινηνού γραφείου}
\begin{enumerate}
  \item Επιλέξτε κύρια App εγγράφων (WPS ή Microsoft 365), login και συνδέστε cloud.
  \item Ανοίξτε τις ση$\mu$ειώσεις, χρησι$\mu$οποιήστε τη λειτουργία σάρωσης για να τραβήξετε ένα έγγραφο, εξάγετε ως PDF και στείλτε το στον εαυτό σας.
  \item Επιλέξτε App cloud, ρυθ$\mu$ίστε αυτό$\mu$ατο αντίγραφο φωτογραφιών ("$\mu$όνο Wi-Fi" για δεδο$\mu$ένα).
  \item Δη$\mu$ιουργήστε $\mu$ια νέα σελίδα στο Notion ή Feishu και δοκι$\mu$άστε να φτιάξετε ένα απλό πρότυπο "Η$\mu$ερολόγιο Εργασίας".
\end{enumerate}

\end{tcbpractice}
